\chapter{斐波那契数列}
记斐波那契生成矩阵
\[
M = \begin{pmatrix} 1 & 1 \\ 1 & 0 \end{pmatrix}.
\]
则对任意整数 \(n\ge 1\)，有
\[
M^n = \begin{pmatrix} F_{n+1} & \mathbf{F_n} \\ \mathbf{F_n} & F_{n-1} \end{pmatrix}
\quad (F_0=0).
\]
后面只取右上角元素来读出关于 \(F_n\) 的关系式。
\section{推导项与项的关系}
\subsection{\texorpdfstring{含$m,n$}{含m,n}}
写出 $M^{m+n} = M^m \cdot M^n$：$$ \begin{pmatrix} \mathbf{F_{m+n+1}} & \mathbf{F_{m+n}} \\ F_{m+n} & F_{m+n-1} \end{pmatrix} = \begin{pmatrix} \mathbf{F_{m+1}} & \mathbf{F_m} \\ F_m & F_{m-1} \end{pmatrix} \begin{pmatrix} \mathbf{F_{n+1}} & \mathbf{F_n} \\ F_n & F_{n-1} \end{pmatrix} $$
看左上角：\[\mathbf{F_{m+n+1} = F_{m+1}F_{n+1} + F_m F_n}\]
看右上角：\[\mathbf{F_{m+n} = F_{m+1}F_n + F_m F_{n-1}}\]
在上式中令 $m=n$，看左上角：\[\mathbf{F_{2n+1} = F_{n+1}^2 + F_n^2}\]
看右上角：\[F_{2n} = F_{n+1}F_n + F_n F_{n-1}\]
把 $F_{n+1}$ 拆成 $F_n + F_{n-1}$，代入得到 \[\mathbf{F_{2n} = F_n^2 + 2F_{n-1}F_n}\]
令 $m=n+r,\ n=n-r$，得
\[F_{n+r}F_{n-r} + F_{n+r-1}F_{n-r-1} = F_{(n+r)+(n-r)-1} = \mathbf{F_{2n-1}}\]
两式相减，得到即 \[F_n^2 - F_{n-r}F_{n+r} = (-1)^{n-r}F_r^2\]
由行列式：\[\det(M) = -1\]利用行列式的乘法性质：$\det(M^n) = (\det M)^n$得：
\[F_{n+1}F_{n-1} - F_n^2=(-1)^n\]
进一步有一般恒等式：
\[F_pF_q - F_uF_v = (-1)^{p+1}F_{u-p}F_{q-u},~~\text{（条件 $p+q = u+v$，且 $p$ 最小）}\] 
因为 $p$ 最小，令 $k = u-p$。根据条件 $p+q = u+v$，必然有 $q = v+k$。左边可以写成 
\[F_p F_{v+k} - F_{p+k} F_v= -\det \begin{pmatrix} F_{p+k} & F_p \\ F_{v+k} & F_v \end{pmatrix} \]
利用结论：$F_{n+k} = F_k F_{n+1} + F_{k-1} F_n$，可将矩阵的第一列拆成两个向量的线性组合：
$$ \begin{pmatrix} F_{p+k} \\ F_{v+k} \end{pmatrix} = F_k \begin{pmatrix} F_{p+1} \\ F_{v+1} \end{pmatrix} + F_{k-1} \begin{pmatrix} F_p \\ F_v \end{pmatrix} $$
后一项的行列式两列成比例，故为 \(0\)，只剩
\[
- F_k \cdot \det \begin{pmatrix} F_{p+1} & F_p \\ F_{v+1} & F_v \end{pmatrix}
= -F_k (F_{p+1}F_v - F_{v+1}F_p).
\]
再利用
\[
F_{p+1}F_v - F_{v+1}F_p = (-1)^p F_{v-p},
\]
便得到
\[
F_pF_q - F_uF_v = (-1)^{p+1}F_{u-p}F_{q-u}.
\]


\section{邻项关系或者隔项关系}
用恒等式 $(a+b)^2+(a-b)^2 = 2a^2+2b^2$，
代入 $$ (F_{n+2} + F_{n+1})^2 + (F_{n+2} - F_{n+1})^2 $$ 得
\[F_{n+3}^2 + F_n^2 = 2F_{n+1}^2 + 2F_{n+2}^2\]
还有一个4次幂综合恒等式：\[F_n^4 = F_{n+2}F_{n+1}F_{n-1}F_{n-2} + 1\]
这只需要调用\[F_n^2 - F_{n-r}F_{n+r} = (-1)^{n-r}F_r^2\]
得到两组公式\[F_{n+1}F_{n-1} = F_n^2 - (-1)^n \cdot 1 = \mathbf{F_n^2 + (-1)^n},~~F_{n+2}F_{n-2} = F_n^2 - (-1)^{n-2} F_2^2 = \mathbf{F_n^2 - (-1)^n}\]
乘起来$$ F_{n+2}F_{n+1}F_{n-1}F_{n-2} = \big(F_n^2 + (-1)^n\big)\big(F_n^2 - (-1)^n\big) = (F_n^2)^2 - \big((-1)^n\big)^2 = \mathbf{F_n^4 - 1} $$

\begin{conclusion}
\[\mathbf{F_{n+k}^2 - F_n^2 = L_k F_n F_{n+k} + F_k^2 (-1)^n} \]
\end{conclusion}
设
\[
M^k = \begin{pmatrix} F_{k+1} & F_k \\ F_k & F_{k-1} \end{pmatrix}.
\]
由 Cayley--Hamilton 定理，
\[
(M^k)^2-(\operatorname{tr}M^k)M^k+(\det M^k)I=0.
\]
再利用
\[
\operatorname{tr}M^k = F_{k+1}+F_{k-1}=L_k,\qquad
\det M^k = F_{k+1}F_{k-1}-F_k^2 = (-1)^k,
\]
可得
\[
(M^k)^2-L_kM^k+(-1)^kI=0.
\]
两边同乘 \(M^n\)，并比较右上角元素，得到
\[
F_{n+2k}=L_kF_{n+k}-(-1)^kF_n.
\]
再将卡塔兰恒等式
\[
F_m^2-F_{m-r}F_{m+r}=(-1)^{m-r}F_r^2
\]
取 \(m=n+k,\ r=k\)，得到
\[
F_{n+k}^2-F_nF_{n+2k}=(-1)^nF_k^2.
\]
代入上一式中的 \(F_{n+2k}\) 后，便得到
\[
F_{n+k}^2-L_kF_nF_{n+k}+(-1)^kF_n^2 = (-1)^nF_k^2.
\]
将两边同除以 \(L_k(F_nF_{n+k})^2\)，再从 \(n=1\) 到无穷求和，可得
\[
\mathbf{ \sum_{n=1}^\infty \frac{1}{F_n F_{n+k}} + \frac{F_k^2}{L_k} \sum_{n=1}^\infty \frac{(-1)^n}{(F_n F_{n+k})^2} = \frac{1}{L_k} \sum_{j=1}^k \frac{1}{F_j^2} } .
\]

当 \(k=1\) 时，\(\frac{F_k^2}{L_k}=1\)，故
\[
\sum_{n=1}^{\infty}\frac{1}{F_n F_{n+1}} + \sum_{n=1}^{\infty}\frac{(-1)^n}{(F_n F_{n+1})^2} = 1.
\]
当 \(k=3\) 时，同样有 \(\frac{F_k^2}{L_k}=1\)，于是
\[
\sum_{n=1}^{\infty}\frac{1}{F_n F_{n+3}} + \sum_{n=1}^{\infty}\frac{(-1)^n}{(F_n F_{n+3})^2} = \frac{9}{16}.
\]
若取 \(k=5\)，则系数变为
\[
\frac{F_5^2}{L_5}=\frac{25}{11},
\]
对应公式为
\[
\mathbf{\sum_{n=1}^\infty \frac{1}{F_n F_{n+5}} + \sum_{n=1}^\infty \frac{25(-1)^n}{11(F_n F_{n+5})^2} = \frac{2161}{9900}}.
\]

原题第三问
\[
\sum \frac{2(-1)^n}{(F_n F_{n+3})^2} - \sum \frac{(-1)^n}{(F_n F_{n+1})^2} = \frac{5}{8}
\]
上式可由前面的恒等式线性组合得到。进一步有
\[
F_a \sum \frac{1}{F_n F_{n+a}} - F_{a+2} \sum \frac{1}{F_n F_{n+a+2}} = \mathbf{\frac{1}{F_{a+1}F_{a+2}}}.
\]
例如把 \(k=3\) 的式子乘以 \(F_3=2\)，把 \(k=5\) 的式子乘以 \(F_5=5\)，再相减，可得到
\[
\mathbf{ \sum_{n=1}^\infty \frac{125(-1)^n}{11(F_n F_{n+5})^2} - \sum_{n=1}^\infty \frac{2(-1)^n}{(F_n F_{n+3})^2} = \frac{131}{3960} } .
\]

\section{各种求和}
矩阵法下，求和可直接化为矩阵等比和。
根据特征方程，有 \[\mathbf{M^2 = M + I}\] （$I$ 为单位矩阵 $\begin{pmatrix} 1 & 0 \\ 0 & 1 \end{pmatrix}$）。
因为
\[
M^2 - M = I \implies M(M-I) = I,
\]
所以
\[
(M-I)^{-1} = M
\]
（同样可得 \((M^2-I)^{-1} = M^{-1}\)）。
于是矩阵和为
\[S_n = \sum_{i=1}^n M^i = M (M^n - I) (M-I)^{-1}= M(M^n - I)\mathbf{M} = \mathbf{M^{n+2} - M^2} \]
比较右上角元素：$M^{n+2}$ 的右上角为 $F_{n+2}$，$M^2$ 的右上角为 $F_2=1$，故
\[\sum_{i=1}^n F_i = F_{n+2} - 1\]
跨次数求和也是一样，求偶数项 $\sum_{i=1}^n M^{2i}$，公比变成 $M^2$：
\[S_{even} = M^2 (M^{2n} - I) (M^2 - I)^{-1} = M^2(M^{2n} - I)\mathbf{M^{-1}} = \mathbf{M^{2n+1} - M}\]
将等号两边的矩阵写开，只看右侧一列：
$$ \begin{pmatrix} \dots & \mathbf{\sum F_{2i}} \\ \dots & \mathbf{\sum F_{2i-1}} \end{pmatrix} = \begin{pmatrix} \dots & \mathbf{F_{2n+1}} \\ \dots & \mathbf{F_{2n}} \end{pmatrix} - \begin{pmatrix} \dots & \mathbf{1} \\ \dots & \mathbf{0} \end{pmatrix} $$
由右上角元素得
\[\mathbf{\sum F_{2i} = F_{2n+1} - 1}\]
由右下角元素得
\[\mathbf{\sum F_{2i-1} = F_{2n}}\]
差分求和$\sum i F_i$的推到就是错位相减：令 $S = \sum_{i=1}^n i M^i$
\[(I-M)S = M + M^2 + \dots + M^n - n M^{n+1} = M^{n+2} - M^2 - n M^{n+1}\]
两边左乘 $(I-M)^{-1}$ 也就是 $\mathbf{-M}$：$$ S = -M(M^{n+2} - M^2 - n M^{n+1}) = \mathbf{-M^{n+3} + M^3 + n M^{n+2}} $$
看右上角元素（注：$F_3=2$）：\[\sum i F_i = -F_{n+3} + 2 + n F_{n+2}=(n+1)F_{n+2} - F_{n+4} + 2\]

\begin{conclusion}
    \[\prod_{k=2}^n \left(1 + \frac{(-1)^{k+1}}{F_k^2}\right) = \frac{F_{n+1}}{F_n}\]
\end{conclusion}
公式：$\prod_{k=2}^n \left(1 + \frac{(-1)^{k+1}}{F_k^2}\right) = \frac{F_{n+1}}{F_n}$
卡西尼恒等式
\[
F_{k+1}F_{k-1} = F_k^2 + (-1)^k
\]
给出
\[
1 + \frac{(-1)^k}{F_k^2} = \frac{F_{k+1}F_{k-1}}{F_k^2}.
\]
因此
\[
\prod_{k=2}^n \left(1 + \frac{(-1)^{k+1}}{F_k^2}\right)
= \prod_{k=2}^n \frac{F_{k-1}F_{k+1}}{F_k^2}
= \frac{F_{n+1}}{F_n}.
\]

\section{整除性质}
若 \(m\mid n\)，设 \(n=km\)。则 \(M^n=(M^m)^k\)。在模 \(F_m\) 的意义下，
$$ M^m = \begin{pmatrix} F_{m+1} & F_m \\ F_m & F_{m-1} \end{pmatrix} \equiv \begin{pmatrix} F_{m+1} & \mathbf{0} \\ \mathbf{0} & F_{m-1} \end{pmatrix} \pmod{F_m} $$
它成为对角矩阵，因此
\[
M^{km} = (M^m)^k \equiv \begin{pmatrix} (F_{m+1})^k & \mathbf{0} \\ \mathbf{0} & (F_{m-1})^k \end{pmatrix} \pmod{F_m}.
\]
比较右上角元素即得
\[
F_n \equiv 0 \pmod{F_m}.
\]

\section{组合数与杨辉三角}
\subsection{杨辉三角}
斐波那契数列的无穷生成函数是 \[\frac{x}{1-x-x^2} = \sum_{n=1}^{\infty} F_n x^n\]
利用等比数列展开 $\frac{1}{1-q}$，令 $q = (x+x^2)$：
$$ \frac{x}{1 - (x + x^2)} = x \sum_{k=0}^{\infty} (x + x^2)^k = \sum_{k=0}^{\infty} x^{k+1} (1 + x)^k $$
利用高中二项式定理把 $(1+x)^k$ 展开为 $\sum_{i=0}^k C_k^i x^i$：
$$\frac{x}{1 - (x + x^2)}  = \sum_{k=0}^{\infty} \sum_{i=0}^k C_k^i x^{k+i+1} $$
为确定 $F_n$，只需提取 $x^n$ 的系数。令指数 $k+i+1 = n$，即 $\mathbf{k = n-1-i}$。将 $k$ 代回后，系数 $C_k^i$ 就变成 $C_{n-1-i}^i$。

\subsection{生成函数}
斐波那契数列的无穷生成函数是 \[\frac{x}{1-x-x^2} = \sum_{n=1}^{\infty} F_n x^n\]
输入 $x = \frac{1}{2}$：右边变成 $\frac{1/2}{1 - 1/2 - 1/4} = 2$。得到
\[ \frac{F_1}{2} + \frac{F_2}{4} + \frac{F_3}{8} + \frac{F_4}{16} + \dots = \mathbf{2}\] 
输入 $x = \frac{1}{10}$：右边变成 $\frac{1/10}{1 - 1/10 - 1/100} = \frac{10}{89}$。\[\frac{F_1}{10} + \frac{F_2}{100} + \frac{F_3}{1000} + \dots = \mathbf{\frac{10}{89}} \]

对
\[
\frac{x}{1-x-x^2} = \sum_{n=1}^{\infty} F_n x^n
\]
两边求导，再同乘 \(x\)，得到
\[
\sum_{k=1}^\infty kF_kx^k = \frac{x(1+x^2)}{(1-x-x^2)^2}.
\]
取 \(x=\frac12\)，便有
$$  \frac{1 \cdot F_1}{2} + \frac{2 \cdot F_2}{4} + \frac{3 \cdot F_3}{8} + \frac{4 \cdot F_4}{16} + \frac{5 \cdot F_5}{32} \dots = \mathbf{10}  $$


\subsection{\texorpdfstring{引入$\pi$和$\e$}{引入π和e}}
\begin{conclusion}
\[\sum_{k=0}^\infty \frac{(-1)^k \cdot \pi^{2k+1}}{(2k+1)!} F_{2k+1} = \mathbf{0} \]
\end{conclusion}
正弦函数的泰勒展开是 $\sin(x) = \sum_{k=0}^\infty \frac{(-1)^k}{(2k+1)!} x^{2k+1}$。将 $x$ 替换成 $\pi M$ 后，得到：
$$ \sin(\pi M) = \sum_{k=0}^\infty \frac{(-1)^k \pi^{2k+1}}{(2k+1)!} M^{2k+1} $$
等式右边的右上角是：\[\sum_{k=0}^\infty \frac{(-1)^k \pi^{2k+1}}{(2k+1)!} F_{2k+1}\]
等式左边 $\sin(\pi M)$ 的右上角可由对角化公式
\[
\frac{f(\alpha)-f(\beta)}{\alpha-\beta}
\]
给出，其中 \(\alpha,\beta\) 为特征方程两根。此处 \(f(x)=\sin(\pi x)\)，且由 \(\alpha+\beta=1\) 得 \(\beta=1-\alpha\)。于是
\[
\sin(\pi\beta)=\sin(\pi-\pi\alpha)=\sin(\pi\alpha),
\]
所以分子为 \(0\)，结论成立。

\begin{conclusion}
\[\sum_{n=0}^\infty \frac{F_n}{n!} = \frac{2\sqrt{e}}{\sqrt{5}} \sinh\left(\frac{\sqrt{5}}{2}\right) \]
\end{conclusion}
由 $e^x = \sum_{n=0}^\infty \frac{x^n}{n!}$ 可得
\[
e^M = \sum_{n=0}^\infty \frac{1}{n!} M^n.
\]
右边矩阵的右上角即为无穷级数 $\sum_{n=0}^\infty \frac{F_n}{n!}$。根据对角化定理，矩阵函数 $f(M)$ 的右上角元素等于
\[
\frac{f(\alpha) - f(\beta)}{\alpha - \beta},
\]
其中 \(\alpha,\beta\) 是黄金比例对应的两个根，分母恰好为 \(\sqrt{5}\)。取 \(f(x)=e^x\)，得
\[
\frac{e^\alpha - e^\beta}{\sqrt{5}}.
\]
再将 \(\alpha,\beta\) 代入并整理：
\[
\frac{1}{\sqrt{5}} \left( e^{\frac{1+\sqrt{5}}{2}} - e^{\frac{1-\sqrt{5}}{2}} \right) = \frac{e^{1/2}}{\sqrt{5}} \left( e^{\frac{\sqrt{5}}{2}} - e^{-\frac{\sqrt{5}}{2}} \right).
\]
括号内正是双曲正弦函数 \(\sinh(x)\) 的标准形式。

\begin{conclusion}
\[\sum_{n=0}^\infty \frac{(-1)^n \cdot \pi^{2n}}{(2n)!} F_{2n} = -\frac{2}{\sqrt{5}} \sin\left(\frac{\sqrt{5}\pi}{2}\right)\]
\end{conclusion}
计算$$ \cos(\pi M) = \sum_{n=0}^\infty \frac{(-1)^n \pi^{2n}}{(2n)!} M^{2n} $$
右边展开的右上角是：\[\sum_{n=0}^\infty \frac{(-1)^n \pi^{2n}}{(2n)!} F_{2n}\]
左边 \(\cos(\pi M)\) 的右上角为
\[
\frac{\cos(\pi \alpha) - \cos(\pi \beta)}{\sqrt{5}}.
\]
由于
\[
\alpha = \frac{1+\sqrt{5}}{2},\qquad \beta = \frac{1-\sqrt{5}}{2},
\]
故
\[
\cos(\pi \alpha) = -\sin\left(\frac{\sqrt{5}\pi}{2}\right),\qquad
\cos(\pi \beta) = \sin\left(\frac{\sqrt{5}\pi}{2}\right).
\]
分子化为
\[
-2\sin\left(\frac{\sqrt{5}\pi}{2}\right),
\]
再除以 \(\sqrt5\) 即得结论。


\subsection{利用二项式展开}
\begin{conclusion}
\[\sum_{k=1}^n \binom{n}{k} F_{k-1} = \mathbf{F_{2n-1} - 1}\]
\end{conclusion}
对特征方程 $M^2 = M+I$ 两边取 $n$ 次方进行二项式展开：
$$ M^{2n} = (M+I)^n = \sum_{k=0}^n \binom{n}{k} M^k $$
取等号两边矩阵的右下角元素。由于 $M^k$ 的右下角恰好是 $F_{k-1}$，左边为 $F_{2n-1}$，右边为 $\sum_{k=0}^n \binom{n}{k} F_{k-1}$。注意完整的二项式求和必须从 $k=0$ 开始。当 $k=0$ 时，这一项是 $\binom{n}{0}F_{-1}$。由 $F_1 = F_0 + F_{-1} \implies 1 = 0 + F_{-1}$，可得 $F_{-1}=1$。因此
\[\sum_{k=1}^n \binom{n}{k} F_{k-1} = \mathbf{F_{2n-1} - 1}\]

\begin{conclusion}
\[F_n + F_{n-1} + \sum_{k=0}^{n-2} F_k 2^{n-2-k} = 2^n\]
\end{conclusion}
该求和可写为
\[
\sum M^k 2^{n-2-k} = 2^{n-2} \sum_{k=0}^{n-2} \left(\frac{M}{2}\right)^k.
\]
利用等比求和公式：$S = 2^{n-2} \cdot [I - (\frac{M}{2})^{n-1}] \cdot (I - \frac{M}{2})^{-1}$。
提取常数后，只需求 $(2I - M)$ 的逆矩阵：
$$ 2I - M = \begin{pmatrix} 2 & 0 \\ 0 & 2 \end{pmatrix} - \begin{pmatrix} 1 & 1 \\ 1 & 0 \end{pmatrix} = \begin{pmatrix} 1 & -1 \\ -1 & 2 \end{pmatrix}, (2I - M)^{-1} = \begin{pmatrix} 2 & 1 \\ 1 & 1 \end{pmatrix}=M^2 $$ 
代回原式：$$ S = 2^{n-2} \cdot \left[I - \frac{M^{n-1}}{2^{n-1}}\right] \cdot 2M^2 = \mathbf{2^{n-1}M^2 - M^{n+1}} $$
比较矩阵右上角可得左边等于 \(2^{n-1}-F_{n+1}\)。再加上前面的 \(F_n+F_{n-1}=F_{n+1}\)，总和即为 \(2^{n-1}\)。

\begin{conclusion}[求$F_{pn}$]
   \[ \mathbf{F_{pn} = \sum_{k=0}^n \binom{n}{k} F_p^k F_{p-1}^{n-k} F_k} \]
\end{conclusion}
由特征方程 \(M^2=M+I\)，可归纳得到
\[
M^p = F_p M + F_{p-1} I.
\]
低次幂写出来便是：
$$ \mathbf{M^2 - M - I = 0 \implies M^2 = M + I} $$
由 \(M^2=M+I\) 出发，继续乘以 \(M\)：
\[M^3 = M \cdot M^2= M(M + I) = \mathbf{M^2} + M= (M + I) + M = \mathbf{2M + 1I}\]
此时系数为 \(2,1\)，对应 \(F_3=2,\ F_2=1\)。
\[M^4 = M \cdot M^3 = M(2M + I) = \mathbf{2M^2} + M= 2(M + I) + M = 2M + 2I + M = \mathbf{3M + 2I}\]
此时系数为 \(3,2\)，对应 \(F_4=3,\ F_3=2\)。
\[M^5 = M \cdot M^4 = M(3M + 2I) = \mathbf{3M^2} + 2M= 3(M + I) + 2M = 3M + 3I + 2M = \mathbf{5M + 3I}\]
于是若
\[
M^k = aM+bI,
\]
则
\[
M^{k+1}=M(aM+bI)=aM^2+bM=(a+b)M+aI.
\]
因此新系数满足与斐波那契数列相同的递推关系，从而得到
\[
\mathbf{M^p = F_p M + F_{p-1} I}.
\]
再对该等式两边取 $n$ 次方，并应用二项式定理：
$$ M^{pn} = (F_p M + F_{p-1} I)^n = \sum_{k=0}^n \binom{n}{k} (F_p M)^k (F_{p-1} I)^{n-k} $$
比较右上角元素，便得到
$$ \mathbf{F_{pn} = \sum_{k=0}^n \binom{n}{k} F_p^k F_{p-1}^{n-k} F_k} $$
输入 $p=3$，（此时 $F_3=2, F_2=1$）：\[\boxed{ F_{3n} = \sum_{k=0}^n \binom{n}{k} 2^k F_k }\]
输入 $p=4$，（此时 $F_4=3, F_3=2$）：\[ \boxed{ F_{4n} = \sum_{k=0}^n \binom{n}{k} 3^k 2^{n-k} F_k }\]

\begin{conclusion}[含$(-1)^n$]
 \[ \sum_{k=0}^n \binom{n}{k} (-1)^k F_{n-2k} = 0 \]
\end{conclusion}
由 \(M^2=M+I\) 可求得 \(M-M^{-1}=I\)。两边取 \(n\) 次方并展开：
\[
\sum_{k=0}^n \binom{n}{k} M^{n-k} (-M^{-1})^k
= \sum_{k=0}^n \binom{n}{k} (-1)^k M^{n-2k}
= I.
\]
比较右上角元素，得到
\[
\boxed{ \sum_{k=0}^n \binom{n}{k} (-1)^k F_{n-2k} = 0 } .
\]
其中负下标按 \(F_{-m} = (-1)^{m+1}F_m\) 理解。

\begin{conclusion}[双项乘积求和]
 \[ \sum_{k=0}^n (F_{k+1}F_{n-k} + F_k F_{n-k-1}) = (n+1)F_n \]
\end{conclusion}
由 \(M^n=M^kM^{n-k}\)，对 \(k=0,1,\dots,n\) 求和得
\[
\sum_{k=0}^n M^kM^{n-k}=(n+1)M^n.
\]
比较右上角元素，即
\[
\begin{pmatrix} F_{k+1} & F_k \end{pmatrix} \cdot \begin{pmatrix} F_{n-k} \\ F_{n-k-1} \end{pmatrix}
= F_{k+1}F_{n-k} + F_k F_{n-k-1},
\]
从而得到
\[ \sum_{k=0}^n (F_{k+1}F_{n-k} + F_k F_{n-k-1}) = (n+1)F_n \]



\begin{conclusion}
    \[F_n^3 + 3 F_{n-3}^3 + F_{n-4}^3 = 3 F_{n-1}^3 + 6 F_{n-2}^3\]
\end{conclusion}
设斐波那契特征方程两根为 \(\alpha,\beta\)。由于 \(F_n\) 可由 \(\alpha^n,\beta^n\) 线性表示，\(F_n^3\) 的递推关系对应的特征根为 \(\alpha^3,\beta^3,-\alpha,-\beta\)。因此可构造特征多项式
\[
P(x) = (x-\alpha^3)(x-\beta^3)(x+\alpha)(x+\beta).
\]
利用 \(\alpha+\beta=1,\ \alpha\beta=-1\) 可化简为
\[
(x^2 - 4x - 1)(x^2 + x - 1) = x^4 - 3x^3 - 6x^2 + 3x + 1.
\]
故数列 \(A_n=F_n^3\) 满足
\[
A_n - 3A_{n-1} - 6A_{n-2} + 3A_{n-3} + A_{n-4} = 0,
\]
即
\[
F_n^3 + 3F_{n-3}^3 + F_{n-4}^3 = 3F_{n-1}^3 + 6F_{n-2}^3.
\]

\section{不等式}
$$ \mathbf{\frac{n-2}{n-1} < \log_{F_{n+1}} F_n < \frac{n-1}{n}} $$
